\chapter{Emergency Quick Reference Cards}\label{ch:quickref}

Use these cards during an emergency. Each card lists the essential actions in order and the critical "do nots." When in doubt, call emergency services.

\section{Unresponsive Person}

\begin{itemize}[leftmargin=*]
\item Shout for help. Tap shoulders and shout, ``Are you OK?''
\item If NO response: call emergency services immediately (send someone else if possible).
\item Open airway (head-tilt/chin-lift). Check normal breathing for 5--10 seconds.
\item \textbf{Not breathing / only gasping:} start CPR (30 compressions, then 2 breaths; repeat).
\item \textbf{Breathing normally:} place in recovery position, monitor continuously.
\item DO NOT leave the person alone.
\end{itemize}

\section{Cardiac Arrest -- CPR}

\begin{itemize}[leftmargin=*]
\item Ensure scene safety. Confirm unresponsive and not breathing normally.
\item Call emergency services. Send for an AED.
\item Push hard and fast: \textbf{100--120 compressions per minute, at least 2 inches (5 cm)} in adults.
\item Allow full chest recoil between compressions; minimize interruptions.
\item 30 compressions : 2 rescue breaths (hands-only CPR if untrained or unwilling).
\item Continue until help arrives, an AED is ready, or the person shows signs of life.
\end{itemize}

\section{Choking -- Conscious Adult and Child (over 1 year)}

\begin{itemize}[leftmargin=*]
\item Ask, ``Are you choking?'' If the person can cough, speak, or breathe, encourage coughing.
\item If the person CANNOT breathe, cough, or speak: give abdominal thrusts (Heimlich).
\item Stand behind, fist above the navel, quick upward-inward thrusts until the object comes out.
\item If the person becomes unconscious: lower to the ground, call emergency services, begin CPR.
\item \textbf{Pregnant or obese:} use chest thrusts instead of abdominal thrusts.
\end{itemize}

\section{Choking -- Conscious Infant (under 1 year)}

\begin{itemize}
\item Support the infant face-down along your forearm, head lower than the chest.
\item Give 5 firm back blows between the shoulder blades.
\item Turn face-up; give 5 chest thrusts with two fingers on the lower half of the sternum.
\item Alternate 5 back blows and 5 chest thrusts until the object is expelled or the infant becomes unconscious.
\item NEVER use abdominal thrusts on an infant.
\end{itemize}

\section{Severe Bleeding}

\begin{itemize}[leftmargin=*]
\item Wear gloves if available. Apply firm, direct pressure with a clean cloth or dressing.
\item Keep pressure for at least 5 minutes without lifting to check.
\item If blood soaks through: add more layers on top -- do NOT remove soaked dressings.
\item Elevate the injured limb above heart level if no fracture is suspected.
\item Life-threatening extremity bleeding that direct pressure fails to stop: apply a tourniquet 2--3 inches above the wound, note the time, and do NOT loosen it.
\item Call emergency services for persistent, heavy, or spurting bleeding.
\end{itemize}

\section{Burns}

\begin{itemize}[leftmargin=*]
\item Remove the person from the source of heat. Ensure your own safety.
\item Cool the burn under cool (not cold) running water for 10--20 minutes.
\item Remove jewelry and tight clothing before swelling begins.
\item Cover with a clean, non-adhesive dressing or plastic film (cling wrap).
\item DO NOT apply ice, butter, oils, or toothpaste. DO NOT break blisters.
\item Seek emergency care for burns larger than a palm, on the face/hands/feet/genitals/joints, or caused by electricity or chemicals.
\end{itemize}

\section{Heart Attack}

\begin{itemize}[leftmargin=*]
\item Have the person stop activity and rest in a comfortable sitting position.
\item Call emergency services immediately. Do not drive them yourself.
\item Assist with their prescribed nitroglycerin as directed.
\item Chew one adult aspirin (325 mg) or four low-dose aspirins (81 mg) if not allergic, not advised against it by a doctor, and a heart attack is likely -- DO NOT give to children or if bleeding is suspected.
\item Monitor breathing. Be ready to begin CPR and use an AED.
\end{itemize}

\section{Stroke -- FAST}

\begin{itemize}[leftmargin=*]
\item \textbf{F -- Face:} Ask to smile. Is one side drooping?
\item \textbf{A -- Arm:} Ask to raise both arms. Does one drift down?
\item \textbf{S -- Speech:} Ask to repeat a phrase. Is speech slurred?
\item \textbf{T -- Time:} Note when symptoms began and call emergency services immediately.
\item Keep the person comfortable; do NOT give food, drink, or medication.
\end{itemize}

\section{Severe Allergic Reaction (Anaphylaxis)}

\begin{itemize}[leftmargin=*]
\item Recognize: hives, swelling of lips/tongue/throat, difficulty breathing, wheezing, dizziness, vomiting, feeling of doom.
\item Give the prescribed epinephrine auto-injector into the outer thigh immediately -- do not wait.
\item Call emergency services immediately.
\item Lie the person flat with legs elevated (or sit them up if they have trouble breathing).
\item A second dose may be given 5--15 minutes later if no improvement and available.
\item Begin CPR if the person becomes unresponsive and stops breathing normally.
\end{itemize}

\section{Seizure}

\begin{itemize}[leftmargin=*]
\item Protect from injury: clear hard objects, pad the head. Do NOT restrain.
\item Turn the person onto their side if possible. Do NOT put anything in the mouth.
\item Time the seizure. Stay until fully recovered.
\item Call emergency services for: first seizure, seizure lasting over 5 minutes, repeated seizures, injury, diabetes/pregnancy, or water.
\end{itemize}

\section{Fainting (Syncope)}

\begin{itemize}[leftmargin=*]
\item Lay the person flat on their back. Loosen tight clothing.
\item Raise the legs only for non-traumatic causes (simple fainting) and only if comfortable.
\item Check breathing and pulse. Begin CPR if not breathing normally.
\item Let them recover slowly -- sit up gradually, do not stand quickly.
\item Call emergency services if the faint happened during exercise, came without warning, or came with palpitations, chest pain, confusion, or injury.
\end{itemize}

\section{Very High Blood Pressure Emergency}

\begin{itemize}[leftmargin=*]
\item Call emergency services immediately if severe headache, vision changes, chest pain, shortness of breath, or confusion occur together or worsen.
\item Sit or lie the person down comfortably and keep them calm.
\item Do NOT give any blood pressure medication that is not theirs.
\item Monitor responsiveness and breathing. Be ready to begin CPR.
\item Do not dismiss the symptoms as "just a headache."
\end{itemize}

\section{Low Blood Sugar (Hypoglycemia)}

\begin{itemize}[leftmargin=*]
\item Conscious: give fast-acting sugar (glucose tablets, juice, regular soda, honey -- about 15 g).
\item Wait 15 minutes; recheck or repeat if still low; follow with a snack or meal.
\item Unconscious or unable to swallow: do NOT give anything by mouth. Apply glucose gel to the inside of the cheek if trained. Call emergency services.
\end{itemize}

\section{Heat Stroke}

\begin{itemize}[leftmargin=*]
\item Signs: body temperature over 104\textdegree F (40\textdegree C), hot/dry skin or heavy sweating, confusion or unconsciousness.
\item Call emergency services immediately.
\item Cool rapidly: cold-water immersion if possible, or pour/spray cool water and fan; ice packs to neck, armpits, groin.
\item Do NOT give fluids to an unconscious person. Do not leave them alone.
\end{itemize}

\section{Hypothermia}

\begin{itemize}[leftmargin=*]
\item Move to shelter, remove wet clothing, wrap in dry blankets.
\item Warm the core first (chest, neck, groin) -- do NOT rub or warm the extremities.
\item Give warm sweet drinks if alert and able to swallow. No alcohol or caffeine.
\item Handle gently. If not breathing normally, begin CPR -- hypothermic patients have been revived after prolonged CPR.
\end{itemize}

\section{Poisoning}

\begin{itemize}[leftmargin=*]
\item Remove the person from the source if safe. Do not expose yourself.
\item Do NOT induce vomiting. Do not give anything by mouth unless told to by poison control.
\item Call poison control (US: \textbf{1-800-222-1222}) or emergency services.
\item Save the container/label and note the substance, amount, and time.
\item Flush chemicals from skin or eyes with copious water for 15--20 minutes.
\end{itemize}

\section{Drowning}

\begin{itemize}[leftmargin=*]
\item Ensure your own safety first. Reach, throw, or row -- do not enter the water unless trained.
\item Call emergency services immediately.
\item Check breathing and pulse. Trained rescuers: not breathing but has a pulse -- give rescue breaths (1 every 5--6 seconds for adults, 1 every 2--3 seconds for child/infant). No pulse, or untrained rescuer: start full CPR (compressions first).
\item Even if the person seems fine, seek medical evaluation for delayed complications.
\end{itemize}

\section{Personal Emergency Numbers}

Fill in before you need them:

\begin{center}
\begin{tabular}{@{}>{\raggedright\arraybackslash}p{5.5cm}p{8cm}@{}}
\hline
Emergency services (112 / 108 in India) & \rule{0pt}{2.5ex}\hrulefill \\
Poison control & \hrulefill \\
Family physician & \hrulefill \\
Local hospital / urgent care & \hrulefill \\
Out-of-area contact & \hrulefill \\
\hline
\end{tabular}
\end{center}

\index{quick reference}
\index{emergency action cards}
\index{CPR quick reference}
\index{emergency numbers}
\index{fainting quick reference}
\index{hypertensive emergency quick reference}
